\documentclass{beamer}
\hypersetup{colorlinks=true, urlcolor=blue}
\usepackage{listings}
\lstdefinestyle{demo}{
    basicstyle=\ttfamily\scriptsize,
    backgroundcolor=\color{black!5},
    frame=single,
    breaklines=true,
    columns=fullflexible,
    keepspaces=true,
}
% Inline code: verbatim argument, so underscores/braces/etc. need no escaping.
% Requires the enclosing frame to be [fragile].
\NewDocumentCommand{\code}{v}{\texttt{#1}}
\title{Python and DataFrames for Sensible Experiment Management}
\author{Ben Keene}
\institute{University of Central Florida}
\date{23 June 2026}

\begin{document}

\frame{\titlepage}

\begin{frame}
    \frametitle{Overview}
    Structured, reproducible workflows for computational research.

    \bigskip

    Today we build a small LLM-benchmarking pipeline end to end, resting on
    one idea:

    \begin{alertblock}{\centering Dataframes as experiment managers:}
        \begin{itemize}
            \item Make parameter sweeps, each run a row.
            \item Save entire sweep as CSV
            \item Load it back in, run each row, write results to disk.
            \item Collate results into a single dataframe for analysis.
        \end{itemize}
    \end{alertblock}
\end{frame}

\begin{frame}
    \frametitle{The Question}
    \begin{itemize}
        \item \url{https://arxiv.org/abs/2205.11916} --- \emph{Large Language
              Models are Zero-Shot Reasoners} (Kojima et al., 2022).
        \item The claim: simply appending
              \textbf{``Let's think step by step''} to a prompt improves
              reasoning.
        \item Our job today: reproduce it. Test it against grade-school arithmetic problems.
    \end{itemize}
\end{frame}

\begin{frame}[fragile]
    \frametitle{Example Problem}
    Think back to 2022, before the explosion of LLMs.

    \bigskip

    Felix has 13 coins. Felix drops 10 coins. Felix drops 2 coins. Felix is
    given 1 more coin. How many coins does Felix have now?

    \bigskip

    This is not an easy question for a 2022 LLM to answer correctly.

    \bigskip

    We have two prompts: a bare \code{standard} prompt and \code{zs_cot}
    (``let's think step by step'').
    We will compare accuracy between these strategies.
\end{frame}

\begin{frame}
    \frametitle{Reproducibility}
    One run proves nothing. To make a claim you'd defend, you need:
    \begin{itemize}
        \item \textbf{Many samples}, different questions.
        \item \textbf{Repeats}, different random seeds.
        \item \textbf{Fixed seeds}, for reproducibility (NB: LLM temperature).
        \item \textbf{Recorded work}, so you can re-run stages of the workflow.
    \end{itemize}

    \bigskip

    That discipline \emph{is} the pipeline. The rest of the talk is its five
    stages.
\end{frame}

\begin{frame}[fragile]
    \frametitle{The Big Picture}
    Each script performs one function.
    \begin{lstlisting}[style=demo]
make_arithmetic.py   -> data/arithmetic.jsonl   the dataset
        |
build_experiment.py  -> experiment.csv     PLAN   1 row / condition
        |
    runner.py        -> runs.csv           RUN    1 row / execution
                         results/{run_id}.jsonl    + per-sample answers
        |
    collate.py       -> collated.csv       COLLATE 1 row / question
        |
        v                plots              ANALYSIS
    \end{lstlisting}
\end{frame}

\begin{frame}[fragile]
    \frametitle{Stage 0: Data Generation}
    \code{make_arithmetic.py} writes \code{data/arithmetic.jsonl}. Each item
    is a dict: a \code{question}, the integer \code{answer}, the \code{steps}
    it took, and an \code{n_steps} count.
    \begin{lstlisting}[style=demo]
>>> dataset[0]
{'question': 'Mia has 12 coins. Mia drops 5 coins. '
             'Mia is given 3 more coins. '
             'How many coins does Mia have now?',
 'answer': 10,
 'n_steps': 2,
 'steps': [-5, 3]}
    \end{lstlisting}
\end{frame}

\begin{frame}[fragile]
    \frametitle{Seeded RNG = Reproducibility}
    The generator seeds Python's RNG from its \code{seed} argument.
    \textbf{What will the output be?}
    \begin{lstlisting}[style=demo]
a = generate_arithmetic(n_steps=3, seed=42)
b = generate_arithmetic(n_steps=3, seed=42)

a == b     # ?
    \end{lstlisting}

    \bigskip

    This is key: When you need to reproduce your plots, you can re-run the pipeline with the same seeds and get the same results.
\end{frame}

\begin{frame}[fragile]
    \frametitle{Stage 1:  PLAN}
    Construct a factory/registry of tasks and LLMs.
    Dict with keys as names, and values are \code{None} because we don't need to execute these tasks yet, only planning right now.
    \begin{lstlisting}[style=demo]
TASK_REGISTRY = {
    "arithmetic": None,
    # "coin_flips": None,
}

LLM_REGISTRY = {
    "qwen2.5-7b": None,
}
    \end{lstlisting}
    \code{build_experiment.py} raises a \code{KeyError} if a sweep references
    a task or LLM that isn't registered.
\end{frame}

\begin{frame}[fragile]
    \frametitle{Stage 1: PLAN, Sweeps}
    Declare an arg space; take its Cartesian product. One \emph{condition} is
    one point in that product.
    \begin{lstlisting}[style=demo]
TASK_SWEEPS = {
    "arithmetic": {
        "condition": ["standard", "zs_cot"],
        "n_samples": [10],
    },
}
LLM_SWEEP = {"llm": ["qwen2.5-7b"],
             "temperature": [0.0], "top_k": [3]}

llm_combos = itertools.product(
    *(LLM_SWEEP[k] for k in keys))
    \end{lstlisting}
\end{frame}

\begin{frame}[fragile]
    \frametitle{Stage 2: RUN, Expansion}
    The plan has one row per \emph{condition}. Running it expands each
    condition into \code{repeats} executions, each with its own \code{run_id}
    and \code{seed}.
    \begin{lstlisting}[style=demo]
for _, cond in plan.iterrows():
    for k in range(int(cond["repeats"])):
        rows.append({
            "run_id": f"run{cond['idx']:02d}_rep{k}",
            "condition_idx": int(cond["idx"]),
            "repeat_idx": k,
            "seed": int(cond["base_seed"]) + k,
        })
    \end{lstlisting}
    That per-run \code{seed} is what makes a single execution reproducible.
\end{frame}

\begin{frame}[fragile]
    \frametitle{Stage 2: RUN, The Live LLM}
    \code{runner.run_all} calls Qwen2.5-7B through \code{llm.py},  an
    OpenAI-compatible client that submits queries to our LLM through a parallelized ThreadPoolExecutor.
    \begin{lstlisting}[style=demo]
model = Qwen(temperature=0.0, top_k=3)

answers = model.chat([
    "Mia has 12 coins. ...",
    "Leo has 8 apples. ...",
])   # -> list of completions, one per prompt
    \end{lstlisting}
    Each run writes its samples to \code{results/{run_id}.jsonl} and a summary
    row to \code{runs.csv}.
\end{frame}

\begin{frame}[fragile]
    \frametitle{Stage 3: Collate}
    \code{collate.py} reads the individual outputs of each run and combines them into a single dataframe.
    \begin{lstlisting}[style=demo]
frames = []
for _, run in runs.iterrows():
    with open(run["results_path"]) as f:
        samples = pd.DataFrame(
            json.loads(line) for line in f)
    for col in RUN_META:
        samples[col] = run[col]
    frames.append(samples)

df = pd.concat(frames, ignore_index=True)
df["correct"] = df["prediction"] == df["answer"]
    \end{lstlisting}
\end{frame}

\begin{frame}
    \frametitle{Getting Started}
    Navigate to the following URL to get started: \url{https://bit.ly/rcijupyter}.

    \bigskip

    Your username is the first letter of your first name, and your full last name.

    $$\text{John Doe} \mapsto \text{jdoe}$$
    $$\text{Ima Knight} \mapsto \text{iknight}$$

    \bigskip

    \begin{alertblock}{\centering Remember your password!}
        Your password will be set to whatever you enter in the password field.
        Make sure to remember it for future use.
    \end{alertblock}
\end{frame}

\begin{frame}[fragile]
    \frametitle{Analysis --- Accuracy by Condition}
    One tidy dataframe, one \code{groupby}. This is the headline:
    \begin{lstlisting}[style=demo]
df.groupby("arg_condition")["correct"].mean()
    \end{lstlisting}

    \bigskip

    \textbf{Predict it first.} Does \code{zs_cot} (``let's think step by
    step'') actually beat the bare \code{standard} prompt?

    \bigskip

    We reveal the real numbers live in the notebook.
\end{frame}

\begin{frame}[fragile]
    \frametitle{Analysis --- Accuracy by Difficulty}
    \code{n_steps} is how many arithmetic operations the problem chains.
    \begin{lstlisting}[style=demo]
df.groupby("n_steps")["correct"].mean()
    \end{lstlisting}

    \bigskip

    \textbf{Predict it.} Do you expect accuracy to fall as the problems get
    longer? By how much?

    \bigskip

    This column only exists because \code{collate} carried it through
    untouched.
\end{frame}

\begin{frame}[fragile]
    \frametitle{Analysis --- The Cost of Reasoning}
    \code{zs_cot} makes the model generate a whole chain of reasoning before
    answering. What does that do to wall-clock time?
    \begin{lstlisting}[style=demo]
(df.groupby(["arg_condition", "run_id"])
   ["runtime_s"].first()
   .groupby("arg_condition").mean())
    \end{lstlisting}

    \bigskip

    The accuracy/cost tradeoff: if \code{zs_cot} wins on accuracy, what did it
    cost you per question?
\end{frame}

\end{document}